% Options for packages loaded elsewhere
\PassOptionsToPackage{unicode}{hyperref}
\PassOptionsToPackage{hyphens}{url}
\documentclass[
  ignorenonframetext,
  brazilian,
  aspectratio=169,
  10pt]{beamer}
\newif\ifbibliography
\usepackage{pgfpages}
\setbeamertemplate{caption}[numbered]
\setbeamertemplate{caption label separator}{: }
\setbeamercolor{caption name}{fg=normal text.fg}
\beamertemplatenavigationsymbolsempty
% remove section numbering
\setbeamertemplate{part page}{
  \centering
  \begin{beamercolorbox}[sep=16pt,center]{part title}
    \usebeamerfont{part title}\insertpart\par
  \end{beamercolorbox}
}
\setbeamertemplate{section page}{
  \centering
  \begin{beamercolorbox}[sep=12pt,center]{section title}
    \usebeamerfont{section title}\insertsection\par
  \end{beamercolorbox}
}
\setbeamertemplate{subsection page}{
  \centering
  \begin{beamercolorbox}[sep=8pt,center]{subsection title}
    \usebeamerfont{subsection title}\insertsubsection\par
  \end{beamercolorbox}
}
% Prevent slide breaks in the middle of a paragraph
\widowpenalties 1 10000
\raggedbottom
\AtBeginPart{
  \frame{\partpage}
}
\AtBeginSection{
  \ifbibliography
  \else
    \frame{\sectionpage}
  \fi
}
\AtBeginSubsection{
  \frame{\subsectionpage}
}
\usepackage{iftex}
\ifPDFTeX
  \usepackage[T1]{fontenc}
  \usepackage[utf8]{inputenc}
  \usepackage{textcomp} % provide euro and other symbols
\else % if luatex or xetex
  \usepackage{unicode-math} % this also loads fontspec
  \defaultfontfeatures{Scale=MatchLowercase}
  \defaultfontfeatures[\rmfamily]{Ligatures=TeX,Scale=1}
\fi
\usepackage{lmodern}
\usetheme[]{Madrid}
\usefonttheme[]{professionalfonts}
\ifPDFTeX\else
  % xetex/luatex font selection
\fi
% Use upquote if available, for straight quotes in verbatim environments
\IfFileExists{upquote.sty}{\usepackage{upquote}}{}
\IfFileExists{microtype.sty}{% use microtype if available
  \usepackage[]{microtype}
  \UseMicrotypeSet[protrusion]{basicmath} % disable protrusion for tt fonts
}{}
\makeatletter
\@ifundefined{KOMAClassName}{% if non-KOMA class
  \IfFileExists{parskip.sty}{%
    \usepackage{parskip}
  }{% else
    \setlength{\parindent}{0pt}
    \setlength{\parskip}{6pt plus 2pt minus 1pt}}
}{% if KOMA class
  \KOMAoptions{parskip=half}}
\makeatother
\usepackage{color}
\usepackage{fancyvrb}
\newcommand{\VerbBar}{|}
\newcommand{\VERB}{\Verb[commandchars=\\\{\}]}
\DefineVerbatimEnvironment{Highlighting}{Verbatim}{commandchars=\\\{\}}
% Add ',fontsize=\small' for more characters per line
\usepackage{framed}
\definecolor{shadecolor}{RGB}{248,248,248}
\newenvironment{Shaded}{\begin{snugshade}}{\end{snugshade}}
\newcommand{\AlertTok}[1]{\textcolor[rgb]{0.94,0.16,0.16}{#1}}
\newcommand{\AnnotationTok}[1]{\textcolor[rgb]{0.56,0.35,0.01}{\textbf{\textit{#1}}}}
\newcommand{\AttributeTok}[1]{\textcolor[rgb]{0.13,0.29,0.53}{#1}}
\newcommand{\BaseNTok}[1]{\textcolor[rgb]{0.00,0.00,0.81}{#1}}
\newcommand{\BuiltInTok}[1]{#1}
\newcommand{\CharTok}[1]{\textcolor[rgb]{0.31,0.60,0.02}{#1}}
\newcommand{\CommentTok}[1]{\textcolor[rgb]{0.56,0.35,0.01}{\textit{#1}}}
\newcommand{\CommentVarTok}[1]{\textcolor[rgb]{0.56,0.35,0.01}{\textbf{\textit{#1}}}}
\newcommand{\ConstantTok}[1]{\textcolor[rgb]{0.56,0.35,0.01}{#1}}
\newcommand{\ControlFlowTok}[1]{\textcolor[rgb]{0.13,0.29,0.53}{\textbf{#1}}}
\newcommand{\DataTypeTok}[1]{\textcolor[rgb]{0.13,0.29,0.53}{#1}}
\newcommand{\DecValTok}[1]{\textcolor[rgb]{0.00,0.00,0.81}{#1}}
\newcommand{\DocumentationTok}[1]{\textcolor[rgb]{0.56,0.35,0.01}{\textbf{\textit{#1}}}}
\newcommand{\ErrorTok}[1]{\textcolor[rgb]{0.64,0.00,0.00}{\textbf{#1}}}
\newcommand{\ExtensionTok}[1]{#1}
\newcommand{\FloatTok}[1]{\textcolor[rgb]{0.00,0.00,0.81}{#1}}
\newcommand{\FunctionTok}[1]{\textcolor[rgb]{0.13,0.29,0.53}{\textbf{#1}}}
\newcommand{\ImportTok}[1]{#1}
\newcommand{\InformationTok}[1]{\textcolor[rgb]{0.56,0.35,0.01}{\textbf{\textit{#1}}}}
\newcommand{\KeywordTok}[1]{\textcolor[rgb]{0.13,0.29,0.53}{\textbf{#1}}}
\newcommand{\NormalTok}[1]{#1}
\newcommand{\OperatorTok}[1]{\textcolor[rgb]{0.81,0.36,0.00}{\textbf{#1}}}
\newcommand{\OtherTok}[1]{\textcolor[rgb]{0.56,0.35,0.01}{#1}}
\newcommand{\PreprocessorTok}[1]{\textcolor[rgb]{0.56,0.35,0.01}{\textit{#1}}}
\newcommand{\RegionMarkerTok}[1]{#1}
\newcommand{\SpecialCharTok}[1]{\textcolor[rgb]{0.81,0.36,0.00}{\textbf{#1}}}
\newcommand{\SpecialStringTok}[1]{\textcolor[rgb]{0.31,0.60,0.02}{#1}}
\newcommand{\StringTok}[1]{\textcolor[rgb]{0.31,0.60,0.02}{#1}}
\newcommand{\VariableTok}[1]{\textcolor[rgb]{0.00,0.00,0.00}{#1}}
\newcommand{\VerbatimStringTok}[1]{\textcolor[rgb]{0.31,0.60,0.02}{#1}}
\newcommand{\WarningTok}[1]{\textcolor[rgb]{0.56,0.35,0.01}{\textbf{\textit{#1}}}}
\ifLuaTeX
\usepackage[bidi=basic,shorthands=off,]{babel}
\else
\usepackage[bidi=default,shorthands=off,]{babel}
\fi
\ifLuaTeX
  \usepackage{selnolig} % disable illegal ligatures
\fi
\setlength{\emergencystretch}{3em} % prevent overfull lines
\providecommand{\tightlist}{%
  \setlength{\itemsep}{0pt}\setlength{\parskip}{0pt}}
\usepackage{amsmath}
\usepackage{amssymb}
\usepackage{booktabs}
\usepackage{xcolor}
\definecolor{AzulAula}{HTML}{1D4ED8}
\definecolor{AzulEscuro}{HTML}{0F172A}
\definecolor{CinzaTexto}{HTML}{334155}
\definecolor{VerdeAula}{HTML}{047857}
\setbeamercolor{structure}{fg=AzulAula}
\setbeamercolor{title}{fg=white,bg=AzulEscuro}
\setbeamercolor{frametitle}{fg=white,bg=AzulEscuro}
\setbeamercolor{normal text}{fg=CinzaTexto,bg=white}
\setbeamercolor{block title}{fg=white,bg=AzulAula}
\setbeamercolor{block body}{fg=AzulEscuro,bg=blue!5}
\setbeamercolor{alertblock title}{fg=white,bg=red!75!black}
\setbeamercolor{alertblock body}{fg=AzulEscuro,bg=red!5}
\setbeamertemplate{navigation symbols}{}
\setbeamertemplate{caption}[numbered]
\setbeamertemplate{footline}[frame number]
\setbeamerfont{caption}{size=\scriptsize}
\setlength{\abovecaptionskip}{2pt}
\setlength{\belowcaptionskip}{0pt}
\renewcommand{\tablename}{Tabela}
\renewcommand{\figurename}{Figura}
\usepackage{bookmark}
\IfFileExists{xurl.sty}{\usepackage{xurl}}{} % add URL line breaks if available
\urlstyle{same}
\hypersetup{
  pdftitle={Aula 8: resíduos e diagnóstico do modelo},
  pdfauthor={FLS 6183 - Métodos Quantitativos de Pesquisa II},
  pdflang={pt-BR},
  hidelinks,
  pdfcreator={LaTeX via pandoc}}

\title{Aula 8: resíduos e diagnóstico do modelo}
\subtitle{CEF, projeção linear, diagnóstico e previsão}
\author{FLS 6183 - Métodos Quantitativos de Pesquisa II}
\date{21 de outubro de 2026}

\begin{document}
\frame{\titlepage}

\begin{frame}{Finalidade do diagnóstico}
\protect\phantomsection\label{finalidade-do-diagnuxf3stico}
Um diagnóstico informa se a análise sustenta o resultado que queremos
produzir.

\begin{table}
\centering
\caption{Objetivo da análise e pergunta de diagnóstico}
\begin{tabular}{p{0.23\textwidth}p{0.66\textwidth}}
\toprule
Objetivo & Pergunta relevante \\
\midrule
Prever $Y$ & As previsões funcionam para novas observações? \\
Resumir uma relação & A projeção linear é um resumo útil? De quais observações depende? \\
Quantificar incerteza & Os erros-padrão acomodam a dispersão e a dependência dos dados? \\
\bottomrule
\end{tabular}
\end{table}

O mesmo padrão nos resíduos pode ter consequências diferentes para cada
objetivo.
\end{frame}

\begin{frame}{Erro populacional e resíduo amostral}
\protect\phantomsection\label{erro-populacional-e-resuxedduo-amostral}
Na representação \(Y_i=\mathbf x_i'\boldsymbol\beta+\varepsilon_i\): \[
\underbrace{\varepsilon_i=Y_i-\mathbf x_i'\boldsymbol\beta}_{\text{erro populacional}},
\qquad
\underbrace{\hat e_i=Y_i-\mathbf x_i'\widehat{\boldsymbol\beta}}_{\text{resíduo amostral}}.
\]

\(\boldsymbol\beta\) é o vetor de parâmetros populacionais;
\(\mathbf x_i'\boldsymbol\beta\) é o valor previsto com esses
parâmetros.

\begin{block}{O que observamos e o que calculamos}
Observamos $Y_i$ e $\mathbf x_i$. Em geral, o parâmetro populacional $\boldsymbol\beta$ é desconhecido, por isso não podemos calcular o erro $\varepsilon_i$.

Depois de estimar $\widehat{\boldsymbol\beta}$, calculamos o resíduo $\hat e_i$ com os dados observados e os coeficientes estimados.
\end{block}

Se \(\boldsymbol\beta\) é conhecido, o erro é calculável. Isso ocorre em
simulações com parâmetros conhecidos e na projeção linear de uma
população finita inteiramente observada em um censo.
\end{frame}

\begin{frame}{Por que o resíduo difere do erro?}
\protect\phantomsection\label{por-que-o-resuxedduo-difere-do-erro}
Substituímos \(Y_i=\mathbf x_i'\boldsymbol\beta+\varepsilon_i\) na
definição do resíduo: \[
\begin{aligned}
\hat e_i
&=Y_i-\mathbf x_i'\widehat{\boldsymbol\beta}\\
&=\mathbf x_i'\boldsymbol\beta+\varepsilon_i-\mathbf x_i'\widehat{\boldsymbol\beta}\\
&=\varepsilon_i-\mathbf x_i'(\widehat{\boldsymbol\beta}-\boldsymbol\beta).
\end{aligned}
\]

\begin{block}{A diferença vem da estimação dos coeficientes}
O resíduo incorpora o termo $-\mathbf x_i'(\widehat{\boldsymbol\beta}-\boldsymbol\beta)$.
Se conhecêssemos $\boldsymbol\beta$, poderíamos calcular o erro. Se $\widehat{\boldsymbol\beta}=\boldsymbol\beta$, resíduos e erros coincidiriam.
\end{block}

A identidade vale também quando \(\boldsymbol\beta\) é o coeficiente do
BLP. Ela não exige uma CEF linear.
\end{frame}

\begin{frame}{Erro no BLP}
\protect\phantomsection\label{erro-no-blp}
A minimização do EQM populacional define o coeficiente e o erro de
projeção: \[
\boldsymbol\beta=\arg\min_{\mathbf b}E[(Y-\mathbf x'\mathbf b)^2],
\qquad
\varepsilon_i=Y_i-\mathbf x_i'\boldsymbol\beta.
\]

Esse erro soma o erro da CEF \(u_i\) à diferença entre a CEF e o BLP: \[
\varepsilon_i=
\underbrace{Y_i-m(\mathbf x_i)}_{u_i}
+\underbrace{m(\mathbf x_i)-\mathbf x_i'\boldsymbol\beta}_{\text{erro de aproximação}}.
\]

\begin{block}{Propriedade do erro de projeção}
$E[\mathbf x_i\varepsilon_i]=\mathbf 0$, mas $E[\varepsilon_i\mid\mathbf x_i]$ pode ser diferente de zero. Se a CEF é linear na especificação escolhida, o erro do BLP coincide com o erro da CEF.
\end{block}

MQO minimiza o critério na amostra, estima \(\boldsymbol\beta\) e produz
os resíduos \(\hat e_i=Y_i-\mathbf x_i'\widehat{\boldsymbol\beta}\).

\scriptsize Hansen, seção 2.18: definição populacional do BLP e do erro
de projeção.
\end{frame}

\begin{frame}{Propriedades por construção}
\protect\phantomsection\label{propriedades-por-construuxe7uxe3o}
Com intercepto e MQO, as equações normais garantem, para cada regressor
incluído: \[
\sum_{i=1}^n\hat e_i=0,
\qquad
\sum_{i=1}^n X_{ij}\hat e_i=0.
\]

Essas igualdades valem mesmo quando a CEF tem outra forma.

\begin{block}{O que é zero por construção}
Regredir os resíduos nos mesmos regressores produz valores ajustados zero.
Isso não garante média condicional zero dos resíduos: uma suavização pode revelar curvatura.
\end{block}

Essas identidades amostrais não testam
\(E[\varepsilon_i\mid\mathbf x_i]=0\).
\end{frame}

\begin{frame}{Curvatura na esperança condicional}
\protect\phantomsection\label{curvatura-na-esperanuxe7a-condicional}
Exemplo sintético: \(X\sim U[-2,2]\), \(u\sim N(0,0{,}6^2)\)
independente de \(X\) e \[
Y=2+1{,}5X-0{,}8(X^2-4/3)+u.
\]

\begin{figure}

{\centering \includegraphics[width=0.74\linewidth]{../output/aula_08/figuras/01_curvatura}

}

\caption{CEF quadrática (verde) e MQO simples estimado na amostra (azul). Fonte: 450 observações sintéticas.}\label{fig:curvatura}
\end{figure}

\small

Aqui, \(E[X^2]=4/3\) e \(\operatorname{Cov}(X,X^2)=0\) dão BLP
\(2+1{,}5X\). Esses coeficientes foram preservados por construção. Em
geral, não basta retirar o termo quadrático da CEF para obter o BLP.
\end{frame}

\begin{frame}{Curvatura nos resíduos}
\protect\phantomsection\label{curvatura-nos-resuxedduos}
Usamos os resíduos da regressão simples do exemplo anterior.

\begin{figure}

{\centering \includegraphics[width=0.82\linewidth]{../output/aula_08/figuras/02_residuos_curvatura} 

}

\caption{Resíduos contra X, com reta de MQO e suavização local (loess). A reta horizontal decorre das equações normais. A curva mostra a estrutura que o modelo simples deixou nos resíduos. Fonte: mesmos 450 dados sintéticos.}\label{fig:residuos-curvatura}
\end{figure}

\begin{block}{Interpretação}
A curvatura indica uma aproximação incompleta da CEF. O coeficiente do BLP continua definido.
\end{block}
\end{frame}

\begin{frame}{Consequências da curvatura}
\protect\phantomsection\label{consequuxeancias-da-curvatura}
\begin{itemize}
\tightlist
\item
  \textbf{Previsão:} o modelo pode errar sistematicamente em regiões de
  \(X\). Uma especificação mais flexível precisa ser avaliada em novos
  dados.
\item
  \textbf{Resumo linear:} o coeficiente continua resumindo a relação na
  população e no suporte considerados.
\end{itemize}

\begin{block}{Linearidade nos coeficientes}
$Y=\alpha+\beta_1X+\beta_2X^2+\varepsilon$ continua sendo uma regressão linear nos coeficientes. Termos quadráticos, splines e interações permitem maior flexibilidade.
\end{block}
\end{frame}

\begin{frame}{Heterocedasticidade}
\protect\phantomsection\label{heterocedasticidade}
Agora a CEF é linear: \(Y=1+1{,}5X+u\), com \(X\sim U[0,4]\),
\(u=(0{,}25+0{,}5X)Z\) e \(Z\sim N(0,1)\) independente de \(X\).

\begin{figure}

{\centering \includegraphics[width=0.82\linewidth]{../output/aula_08/figuras/03_heterocedasticidade} 

}

\caption{Resíduos em 600 observações sintéticas. A variância do erro aumenta com X; a CEF permanece linear. A linha verde é uma suavização local e a azul marca zero. Fonte: processo gerador descrito.}\label{fig:heterocedasticidade}
\end{figure}

\small

A mudança da dispersão e a curvatura da média são aspectos diferentes do
diagnóstico.
\end{frame}

\begin{frame}{Heterocedasticidade: o caso geral}
\protect\phantomsection\label{heterocedasticidade-o-caso-geral}
A formulação geral permite que a variância do erro dependa dos
regressores: \[
\operatorname{Var}(\varepsilon_i\mid\mathbf x_i)=\sigma^2(\mathbf x_i).
\]

A homocedasticidade é o caso particular em que
\(\sigma^2(\mathbf x_i)=\sigma^2\) para todos os valores dos
regressores.

\begin{block}{Consequência para os erros-padrão}
Os erros-padrão convencionais usam a hipótese de variância constante.
Com heterocedasticidade, podem subestimar ou superestimar a incerteza dos coeficientes.
\end{block}

\begin{itemize}
\tightlist
\item
  Erros-padrão robustos à heterocedasticidade acomodam variâncias
  diferentes entre observações.
\item
  Com observações independentes e condições usuais para amostras
  grandes, podem ser usados tanto com heterocedasticidade quanto com
  homocedasticidade.
\end{itemize}

\scriptsize Hansen, seções 4.13--4.15. Na aula 9: comparação entre
erros-padrão convencionais, HC3 e bootstrap.
\end{frame}

\begin{frame}{Da estimativa de MQO à matriz hat}
\protect\phantomsection\label{da-estimativa-de-mqo-uxe0-matriz-hat}
Partimos da fórmula de MQO, com \(\mathbf X'\mathbf X\) inversível: \[
\widehat{\boldsymbol\beta}=(\mathbf X'\mathbf X)^{-1}\mathbf X'\mathbf Y.
\]

Substituímos
\(\mathbf Y=\mathbf X\boldsymbol\beta+\boldsymbol\varepsilon\): \[
\begin{aligned}
\widehat{\boldsymbol\beta}
&=(\mathbf X'\mathbf X)^{-1}\mathbf X'(\mathbf X\boldsymbol\beta+\boldsymbol\varepsilon)\\
&=\boldsymbol\beta+(\mathbf X'\mathbf X)^{-1}\mathbf X'\boldsymbol\varepsilon.
\end{aligned}
\]

Multiplicando por \(\mathbf X\), obtemos os valores ajustados: \[
\widehat{\mathbf Y}=\mathbf X\widehat{\boldsymbol\beta}
=\mathbf X\boldsymbol\beta+
\underbrace{\mathbf X(\mathbf X'\mathbf X)^{-1}\mathbf X'}_{\mathbf H}
\boldsymbol\varepsilon.
\]

A matriz \(\mathbf H\) leva \(\mathbf Y\) a \(\widehat{\mathbf Y}\):
\(\widehat{\mathbf Y}=\mathbf H\mathbf Y\). Por isso é chamada de matriz
\textit{hat} ou matriz de projeção.
\end{frame}

\begin{frame}{Alavancagem e variância dos resíduos}
\protect\phantomsection\label{alavancagem-e-variuxe2ncia-dos-resuxedduos}
Defina a matriz de projeção e sua diagonal: \[
\mathbf H=\mathbf X(\mathbf X'\mathbf X)^{-1}\mathbf X',
\qquad h_{ii}=\mathbf x_i'(\mathbf X'\mathbf X)^{-1}\mathbf x_i.
\]

\(h_{ii}\) mede o quanto a observação é incomum no conjunto dos
regressores. Com \(p=k+1\) coeficientes, \(0\le h_{ii}\le1\) e
\(\sum_i h_{ii}=p\).

\begin{block}{Erros homocedásticos e não correlacionados}
\[
\operatorname{Var}(\boldsymbol\varepsilon\mid\mathbf X)=\sigma^2\mathbf I
\quad\Longrightarrow\quad
\operatorname{Var}(\hat e_i\mid\mathbf X)=(1-h_{ii})\sigma^2.
\]
Mesmo com erros homocedásticos, os resíduos brutos podem ter variâncias diferentes.
\end{block}

Para \(h_{ii}<1\), dividir \(\hat e_i\) por \(\sqrt{1-h_{ii}}\) corrige
essa diferença de escala sob essas condições.

\scriptsize Hansen, seções 3.19 e 4.10. Alavancagem depende de todos os
regressores, inclusive indicadoras.
\end{frame}

\begin{frame}{Alavancagem e influência}
\protect\phantomsection\label{alavancagem-e-influuxeancia}
Cada painel acrescenta um ponto aos mesmos 35 dados de base.

\begin{figure}

{\centering \includegraphics[width=0.82\linewidth]{../output/aula_08/figuras/04_alavancagem_influencia} 

}

\caption{Triângulo: observação acrescentada. Linha contínua: ajuste com o ponto; tracejada: ajuste sem ele. Um resíduo grande e uma alavancagem alta podem ocorrer separadamente. Fonte: dados sintéticos, mesmos eixos nos painéis.}\label{fig:influencia}
\end{figure}

\small

Uma observação é influente para um resultado quando sua retirada o
altera de forma substantivamente relevante.
\end{frame}

\begin{frame}{Sensibilidade a uma observação}
\protect\phantomsection\label{sensibilidade-a-uma-observauxe7uxe3o}
A inclinação é medida em unidades de \(Y\) por unidade de \(X\).

\begin{table}

\caption{\label{tab:tabela-influencia}Sensibilidade da inclinação nos exemplos sintéticos}
\centering
\begin{tabular}[t]{lrrr}
\toprule
Ponto acrescentado & $h_{ii}$ & Resíduo & $\hat\beta_1-\hat\beta_{1(-i)}$\\
\midrule
Resíduo grande & 0,028 & 5,809 & 0,000\\
Alavancagem alta & 0,481 & 0,000 & 0,000\\
Ambos & 0,481 & −3,495 & −0,629\\
\bottomrule
\end{tabular}
\end{table}

\begin{block}{O que investigar}
Conferir registro e unidade de análise; compreender a observação; comparar os resultados com e sem ela e justificar a decisão substantiva.
\end{block}

Uma observação incomum pode conter informação relevante. Um limiar de
diagnóstico não determina sua exclusão.

\scriptsize Hansen, seções 3.20--3.21. A influência depende do
coeficiente ou previsão de interesse.
\end{frame}

\begin{frame}{Dependência na ordem dos dados}
\protect\phantomsection\label{dependuxeancia-na-ordem-dos-dados}
Exemplo sintético: \(Y_t=2+1{,}5X_t+u_t\), com \(X_t\) independente do
erro e \(u_t=0{,}8u_{t-1}+v_t\), \(v_t\sim N(0,1)\) independente ao
longo do tempo.

\begin{figure}

{\centering \includegraphics[width=0.82\linewidth]{../output/aula_08/figuras/05_dependencia} 

}

\caption{Mesmos 220 resíduos na ordem temporal e em uma ordem aleatória. Reordenar as linhas mantém o ajuste de MQO e pode ocultar a estrutura serial. Fonte: dados sintéticos.}\label{fig:dependencia}
\end{figure}

\small

Diagnósticos temporais e espaciais exigem uma ordem ou distância com
significado. Um identificador numérico não constitui uma coordenada
geográfica.
\end{frame}

\begin{frame}{Dependência dos resíduos: exemplo com intercepto}
\protect\phantomsection\label{dependuxeancia-dos-resuxedduos-exemplo-com-intercepto}
Considere uma regressão apenas com intercepto, com erros independentes
de média zero: \[
Y_i=\mu+\varepsilon_i.
\]

MQO estima o intercepto pela média amostral: \[
\widehat\mu=\bar Y=\mu+\bar\varepsilon.
\]

Substituindo na definição do resíduo: \[
\begin{aligned}
\hat e_i&=Y_i-\bar Y\\
&=(\mu+\varepsilon_i)-(\mu+\bar\varepsilon)\\
&=\varepsilon_i-\bar\varepsilon.
\end{aligned}
\]

\begin{block}{A dependência é criada pelo ajuste}
Todos os resíduos compartilham a mesma $\bar\varepsilon$ e somam zero.
Com vários regressores, $\widehat{\mathbf e}=(\mathbf I-\mathbf H)\boldsymbol\varepsilon$.
Essa dependência, por si só, não exige erros-padrão agrupados.
\end{block}
\end{frame}

\begin{frame}[fragile]{Dependência nos erros e inferência}
\protect\phantomsection\label{dependuxeancia-nos-erros-e-inferuxeancia}
No exemplo \(u_t=0{,}8u_{t-1}+v_t\), um choque pode persistir por vários
períodos. HC3 acomoda heterocedasticidade, mas não essa dependência
serial.

\begin{block}{Série temporal: erros-padrão HAC, como Newey--West}
HAC significa consistentes sob heterocedasticidade e autocorrelação.
Eles incorporam covariâncias entre períodos, sob condições de dependência fraca e amostras grandes.
O ajuste de MQO é mantido; mudam os erros-padrão e os intervalos.
\end{block}

\begin{itemize}
\tightlist
\item
  \textbf{Ordem e defasagens:} ordene os dados pelo tempo. Escolha
  quantas defasagens considerar e examine a sensibilidade a essa
  escolha.
\item
  \textbf{Outras estruturas:} com dependência dentro de grupos, use
  erros-padrão agrupados. Na reamostragem, use grupos inteiros ou blocos
  de períodos consecutivos.
\end{itemize}

Em R, para um modelo ajustado com os dados em ordem temporal, com
seleção automática de defasagens:

\begin{Shaded}
\begin{Highlighting}[]
\FunctionTok{sqrt}\NormalTok{(}\FunctionTok{diag}\NormalTok{(sandwich}\SpecialCharTok{::}\FunctionTok{NeweyWest}\NormalTok{(m)))  }\CommentTok{\# erros{-}padrão}
\end{Highlighting}
\end{Shaded}

\scriptsize Fonte: documentação de
\href{https://sandwich.r-forge.r-project.org/reference/NeweyWest.html}{NeweyWest}
e
\href{https://sandwich.r-forge.r-project.org/reference/vcovCL.html}{vcovCL},
pacote \texttt{sandwich}.
\end{frame}

\begin{frame}{Aplicação: eleições de 2022}
\protect\phantomsection\label{aplicauxe7uxe3o-eleiuxe7uxf5es-de-2022}
Retomamos a base da aula 7. Uma linha representa uma zona eleitoral.

\begin{block}{Variáveis e comparação}
$Y_i$: porcentagem de Lula no segundo turno.
$X_i$: porcentagem de Lula no primeiro turno.
Cada percentual usa todos os votos válidos do respectivo turno.
\end{block}

Comparamos os modelos \(Y_i=\alpha+\beta_1X_i+\varepsilon_i\) e \[
Y_i=\alpha+\beta_1X_i+\sum_{g\ne AC}\gamma_g\,1\{UF_i=g\}+\varepsilon_i.
\]

Os gráficos usam as mesmas 2.636 zonas das 27 UFs, com peso igual no
ajuste. A observação agregada do Exterior (ZZ) será examinada
separadamente.

\small

Objetivo descritivo e preditivo.

\scriptsize Fonte: dados TSE preservados e documentados para o
laboratório da aula 7.
\end{frame}

\begin{frame}{Resíduos e valores ajustados}
\protect\phantomsection\label{resuxedduos-e-valores-ajustados}
\begin{figure}

{\centering \includegraphics[width=0.82\linewidth]{../output/aula_08/figuras/08_eleitoral_residuos} 

}

\caption{Resíduos dos dois modelos nas mesmas 2.636 zonas das 27 UFs, eleição de 2022. Azul: zero; verde: suavização local. Eixos comuns. Resíduos em pontos percentuais. Fonte: TSE, base preservada do curso.}\label{fig:eleitoral-residuos}
\end{figure}

\begin{block}{Leitura do gráfico}
Onde há regiões com resíduos predominantemente positivos ou negativos?
Como a inclusão de UF altera a estrutura visível e a dispersão?
\end{block}
\end{frame}

\begin{frame}{Magnitude dos resíduos}
\protect\phantomsection\label{magnitude-dos-resuxedduos}
\begin{figure}

{\centering \includegraphics[width=0.82\linewidth]{../output/aula_08/figuras/09_eleitoral_dispersao} 

}

\caption{Valor absoluto dos resíduos contra os valores ajustados, nas mesmas zonas e modelos da figura anterior. Verde: suavização local. Eixos comuns; magnitudes em pontos percentuais. Fonte: TSE, eleição de 2022.}\label{fig:eleitoral-dispersao}
\end{figure}

\small

Uma mudança na magnitude dos resíduos pode refletir variância
condicional variável e também erro de aproximação da média. A leitura
conjunta dos dois gráficos ajuda a distinguir essas possibilidades.
\end{frame}

\begin{frame}{Influência na inclinação eleitoral}
\protect\phantomsection\label{influuxeancia-na-inclinauxe7uxe3o-eleitoral}
Modelo com UF. Calculamos a inclinação novamente ao excluir cada zona.

\small
\begin{table}

\caption{\label{tab:tabela-eleitoral-influencia}Quatro maiores mudanças absolutas na inclinação ao excluir uma zona}
\centering
\begin{tabular}[t]{lrrrr}
\toprule
UF & Zona & $h_{ii}$ & Resíduo (p.p.) & $\Delta\hat\beta_1$\\
\midrule
CE & 121 & 0,010 & 8,56 & −0,00048\\
PA & 91 & 0,016 & −2,46 & 0,00040\\
CE & 24 & 0,010 & 8,80 & −0,00039\\
SP & 1 & 0,004 & 3,84 & 0,00029\\
\bottomrule
\end{tabular}
\end{table}

\(\Delta\hat\beta_1=\hat\beta_1-\hat\beta_{1(-i)}\). Inclinação com
todas as zonas: \(1{,}0111\) p.p. de \(Y\) por 1 p.p. de \(X\).

A relevância da mudança depende do uso do modelo. Uma zona pode ser
importante para uma previsão específica e pouco importante para essa
inclinação comum.
\end{frame}

\begin{frame}{Exterior: ajuste perfeito de uma observação}
\protect\phantomsection\label{exterior-ajuste-perfeito-de-uma-observauxe7uxe3o}
Incluindo ZZ na regressão com categorias geográficas, há 2.637 linhas.
Uma única observação pertence à categoria Exterior.

\begin{block}{Consequência da indicadora exclusiva}
O intercepto da categoria ajusta exatamente sua observação:
$\hat e_{ZZ}\simeq0$ e $h_{ZZ}=1$.
A variância do erro nessa categoria não pode ser estimada a partir de variação interna.
\end{block}

\begin{itemize}
\tightlist
\item
  O resíduo zero decorre da flexibilidade do ajuste para essa
  observação.
\item
  As correções por \(1-h_{ii}\) ficam indefinidas nesse ponto; HC2, HC3
  e exclusão da observação exigem cuidado.
\item
  Nos gráficos anteriores usamos as 27 UFs para comparar modelos na
  mesma amostra. ZZ permanece preservada e explicitamente discutida.
\end{itemize}

\scriptsize Hansen, seção 4.16, sobre indicadoras raras. Ajuste perfeito
não garante precisão para a previsão de outro ano.
\end{frame}

\begin{frame}{Avaliação de novas zonas da mesma eleição}
\protect\phantomsection\label{avaliauxe7uxe3o-de-novas-zonas-da-mesma-eleiuxe7uxe3o}
Divisão única, aleatória dentro de UF: 2121 zonas de treino e 515 de
teste. As especificações foram fixadas antes da comparação.

\small
\begin{table}

\caption{\label{tab:tabela-previsao}Erros de previsão em pontos percentuais, eleição de 2022}
\centering
\begin{tabular}[t]{lrrr}
\toprule
Modelo & REQM treino & REQM teste & EAM teste\\
\midrule
Primeiro turno & 1,51 & 1,39 & 1,04\\
Primeiro turno + UF & 1,20 & 1,14 & 0,85\\
Primeiro turno + quadrado + UF & 1,20 & 1,14 & 0,85\\
\bottomrule
\end{tabular}
\end{table}

REQM: raiz do erro quadrático médio. EAM: erro absoluto médio. Cada zona
tem peso igual.

Esta divisão avalia zonas reservadas de 2022. Uma única divisão tem
variabilidade e pode aproveitar semelhanças entre zonas. Ela não avalia
dependência espacial nem demonstra capacidade de transportar o modelo
para 2026.
\end{frame}

\begin{frame}{Diagnóstico e consequência para a análise}
\protect\phantomsection\label{diagnuxf3stico-e-consequuxeancia-para-a-anuxe1lise}
\small
\begin{table}
\centering
\caption{Aspectos do diagnóstico e decisões pertinentes}
\begin{tabular}{p{0.26\textwidth}p{0.64\textwidth}}
\toprule
Aspecto & Consequência a avaliar \\
\midrule
Curvatura & Aproximação da CEF e previsão fora da amostra. \\
Dispersão variável & Estimação da variância e intervalos compatíveis com heterocedasticidade. \\
Dependência & Unidade de informação, grupos e ordem temporal ou espacial. \\
Influência & Qualidade dos registros e sensibilidade do resultado de interesse. \\
\bottomrule
\end{tabular}
\end{table}

Um padrão merece investigação quando pode comprometer o objetivo ou
delimitar a interpretação do resultado.
\end{frame}

\begin{frame}[fragile]{Exploração no R}
\protect\phantomsection\label{explorauxe7uxe3o-no-r}
O mesmo modelo permite examinar média dos resíduos, magnitude e
alavancagem.

\begin{Shaded}
\begin{Highlighting}[]
\NormalTok{m }\OtherTok{\textless{}{-}} \FunctionTok{lm}\NormalTok{(pct\_lula\_t2 }\SpecialCharTok{\textasciitilde{}}\NormalTok{ pct\_lula\_t1 }\SpecialCharTok{+}\NormalTok{ uf, }\AttributeTok{data =}\NormalTok{ brasil)}
\FunctionTok{plot}\NormalTok{(}\FunctionTok{fitted}\NormalTok{(m), }\FunctionTok{resid}\NormalTok{(m), }\AttributeTok{xlab =} \StringTok{"Valor ajustado (\%)"}\NormalTok{,}
     \AttributeTok{ylab =} \StringTok{"Resíduo (p.p.)"}\NormalTok{)}
\FunctionTok{abline}\NormalTok{(}\AttributeTok{h =} \DecValTok{0}\NormalTok{, }\AttributeTok{col =} \StringTok{"blue"}\NormalTok{)}
\FunctionTok{lines}\NormalTok{(}\FunctionTok{lowess}\NormalTok{(}\FunctionTok{fitted}\NormalTok{(m), }\FunctionTok{resid}\NormalTok{(m)), }\AttributeTok{col =} \StringTok{"darkgreen"}\NormalTok{)}
\FunctionTok{head}\NormalTok{(}\FunctionTok{sort}\NormalTok{(}\FunctionTok{hatvalues}\NormalTok{(m), }\AttributeTok{decreasing =} \ConstantTok{TRUE}\NormalTok{))}
\end{Highlighting}
\end{Shaded}

Para uma zona relevante, compare o coeficiente ou a previsão com e sem
sua observação. Registre o diagnóstico, sua consequência e a
justificativa da decisão.
\end{frame}

\begin{frame}{Leituras}
\protect\phantomsection\label{leituras}
\begin{block}{Leituras previstas no programa}
Hansen, Bruce E. \textit{Econometrics}. Seções 4.10--4.15:
resíduos, variância do erro, previsão e erros-padrão.

Shalizi, Cosma R. \textit{The Truth about Linear Regression}.
Seções 6.1 e 6.5: diagnósticos dos resíduos e projeção sob especificação incorreta.
\end{block}

\small

Complementos: Hansen, seções 3.19--3.21 (alavancagem e influência) e
4.16 (indicadoras raras); Galdino, \textit{Regressão}, capítulo
``Checagem''.

Na aula 9: distribuição amostral, erros-padrão convencionais e robustos
e bootstrap.
\end{frame}

\end{document}
